\begin{proofbox}
	\n\textbf{\textcolor{CProof}{思路提示}}\n\n通过构造函数，利用导数求出函数的最大值或最小值，与零比较。\n\n

	\medskip
	\n\textbf{\textcolor{CProof}{正式推导}}\n
	\begin{description}[style=nextline, leftmargin=2.8em, labelsep=0.5em, itemsep=0.55em]
		\n
		\item[\textbf{步骤一（构造函数对数）：}]
			令 $f(x)=\ln x-(x-1)$，定义域为 $x>0$。\n
			\[
				\n f(x)=\ln x-x+1.\n
			\]
			\n
			\par\textbf{理由：}\n 欲证 $\ln x\le x-1$，等价于证 $f(x)\le 0$，故研究 $f(x)$ 的最大值。\n\n
		\item[\textbf{步骤二（求导对数）：}]
			计算 $f'(x)=\frac{1}{x}-1$，即 $f'(x)=\frac{1-x}{x}$。\n
			\par\textbf{理由：}\n 通过导函数符号判断单调性。\n\n
		\item[\textbf{步骤三（单调性与最值对数）：}]
			当 $0<x<1$ 时 $f'(x)>0$，$f(x)$ 递增；当 $x>1$ 时 $f'(x)<0$，$f(x)$ 递减。所以最大值 $f(1)=0$。\n
			\[
				\n f(x)\le 0\;\Longrightarrow\;\ln x\le x-1.\n
			\]
			\n
			\par\textbf{理由：}\n 最大值点为 $x=1$，代入得 $f(1)=\ln1-1+1=0$。\n\n
		\item[\textbf{步骤四（构造函数指数）：}]
			令 $g(x)=e^x-(x+1)$，定义域为 $\mathbb{R}$。\n
			\[
				\n g(x)=e^x-x-1.\n
			\]
			\n
			\par\textbf{理由：}\n 欲证 $e^x\ge x+1$，等价于证 $g(x)\ge 0$，研究 $g(x)$ 的最小值。\n\n
		\item[\textbf{步骤五（求导指数）：}]
			计算 $g'(x)=e^x-1$。\n
			\par\textbf{理由：}\n 导函数符号决定单调性。\n\n
		\item[\textbf{步骤六（单调性与最值指数）：}]
			当 $x<0$ 时 $g'(x)<0$，$g(x)$ 递减；当 $x>0$ 时 $g'(x)>0$，$g(x)$ 递增。所以最小值 $g(0)=0$。\n
			\[
				\n g(x)\ge 0\;\Longrightarrow\;e^x\ge x+1.\n
			\]
			\n
			\par\textbf{理由：}\n 最小值为 $g(0)=e^0-0-1=0$。\n\n
	\end{description}
	\n\n

	\medskip
	\n\textbf{\textcolor{CProof}{结论回扣}}\n\n通过构造辅助函数并利用导数，我们分别证明了对数不等式 $\ln x\le x-1$（$x>0$）和指数不等式 $e^x\ge x+1$（$x\in\mathbb{R}$），且等号成立条件也由极值点自然给出。这两个不等式在导数放缩中可以直接使用。\n
\end{proofbox}
